\documentclass[11pt,a4paper]{article}
\input{casper_preambule}
\begin{document}
\begin{center}
{\large\bfseries Préparation au Concours National Commun --- Informatique MP}\\[8pt]
\rule{0.9\linewidth}{0.5pt}\\[6pt]
{\Large\bfseries Épreuve d'Informatique}\\[4pt]
{\bfseries Filière MP --- Durée : 4 heures}\\[6pt]
{\itshape Finalité dans les chaînes de blocs : le protocole Casper}\\[2pt]
{\small d'après V.~Buterin et V.~Griffith, \emph{Casper the Friendly Finality Gadget}, \\ Ethereum Foundation, arXiv:1710.09437v4 (2019)}\\[6pt]
\rule{0.9\linewidth}{0.5pt}\\[8pt]
{\small Auteur~: \href{https://orcid.org/0000-0002-6108-7176}{\textbf{Pr Agrégé El Hadiq Zouhair}}}\\[3pt]
{\small \href{https://cpgeacademy.org}{\textbf{cpgeacademy.org}} \quad---\quad Tous droits réservés \textcopyright\ cpgeacademy.org}
\end{center}
\medskip
\begin{center}
\fcolorbox{gray}{gray!5}{\begin{minipage}{0.93\linewidth}\small
\textbf{Avis important.} Ce document est une \textbf{initiative personnelle} du professeur,
à but strictement pédagogique. Il ne s'agit \textbf{pas} d'un sujet officiel du Concours
National Commun~; il s'en inspire uniquement par la forme et l'esprit.
\smallskip

\textbf{Pour aller plus loin.} Les élèves peuvent traiter une \textbf{nouvelle version} de
ce problème dotée d'un \textbf{autocorrecteur des réponses}, et suivre un \textbf{cours
complet}, sur \href{https://cpgeacademy.org}{\textbf{cpgeacademy.org}}.
\end{minipage}}
\end{center}
\medskip


\section*{Présentation générale}

Une \emph{chaîne de blocs} (\emph{blockchain}) est une structure de données répartie dans
laquelle des blocs sont ajoutés les uns à la suite des autres. En cas de latence du réseau
ou d'attaque, plusieurs blocs peuvent être proposés comme successeurs d'un même bloc~: la
structure obtenue n'est plus une liste mais un \textbf{arbre de blocs}, dont la racine est
le \emph{bloc génésis}.

Le protocole \textbf{Casper} (Buterin et Griffith, 2019) est une surcouche de
\emph{finalisation} qui sélectionne, dans cet arbre, une unique chaîne canonique. Pour des
raisons d'efficacité, Casper ne raisonne pas sur l'arbre complet mais sur le sous-arbre des
\textbf{points de contrôle} (\emph{checkpoints})~: le bloc génésis est un point de contrôle,
et tout bloc dont le numéro dans l'arbre est un multiple exact de $100$ en est un également.

Un ensemble de \textbf{validateurs} possédant chacun un \emph{dépôt} (une somme immobilisée)
émettent des \textbf{votes}. Lorsqu'une proportion suffisante du dépôt total soutient un
même vote, un \emph{lien de supermajorité} est créé~; de proche en proche, des points de
contrôle deviennent \emph{justifiés} puis \emph{finalisés}. Deux règles, appelées
\textbf{commandements de Casper}, interdisent certains votes~; tout validateur qui les
enfreint perd la totalité de son dépôt (\emph{slashing}).

Ce problème étudie, de manière algorithmique et mathématique, les objets et les propriétés
de ce protocole. Il comporte \textbf{cinq parties largement indépendantes}, de difficulté
croissante. Les fonctions écrites dans une question peuvent être librement réutilisées dans
les questions suivantes.

\subsection*{Notations et conventions}

\begin{itemize}[leftmargin=1.4em,itemsep=2pt]
\item L'arbre des points de contrôle possède $n$ n\oe uds numérotés de $0$ à $n-1$. Le
n\oe ud $0$ est la \textbf{racine} $r$ (le bloc génésis).
\item L'arbre est représenté par une liste Python \code{parent} de longueur $n$ telle que
\code{parent[0] == -1} et, pour tout $c\ge 1$, \code{parent[c]} est le numéro du père de
$c$. On suppose toujours~: \[\forall c\ge 1,\qquad 0\le \texttt{parent[c]} < c .\]
\item La \textbf{hauteur} $h(c)$ d'un point de contrôle $c$ est le nombre d'arêtes du
chemin de $c$ à la racine. Ainsi $h(r)=0$.
\item Un point de contrôle $a$ est un \textbf{ancêtre} de $b$ si $a$ appartient au chemin
de $b$ à la racine (donc $b$ est ancêtre de lui-même). Deux points de contrôle $a$ et $b$
sont dits \textbf{en conflit} lorsque aucun des deux n'est ancêtre de l'autre.
\item Les validateurs sont identifiés par des clés. Les dépôts sont donnés par un
dictionnaire Python \code{depots} associant à chaque validateur un entier
\emph{strictement positif}. On note $T=\sum_{\nu} d_\nu$ le
\textbf{dépôt total} et, pour un ensemble $E$ de validateurs, $w(E)=\sum_{\nu\in E} d_\nu$.
\item Un ensemble $E$ de validateurs est une \textbf{supermajorité} lorsque
$w(E)\ \ge\ \tfrac{2}{3}\,T$.
\item Un \textbf{vote} est un triplet Python \code{(v, s, t)} où \code{v} est un
validateur, \code{s} la \emph{source} et \code{t} la \emph{cible}. Le vote est
\textbf{valide} si \code{v} appartient au jeu de validateurs et si \code{s} est un ancêtre
\emph{strict} de \code{t}. Sauf mention contraire, tous les votes manipulés sont valides~;
en particulier $h(s)<h(t)$.
\item En Python, seules les constructions du programme officiel sont autorisées~: types de
base, listes, tuples, dictionnaires, ensembles, boucles, fonctions, récursivité, ainsi que
les modules \code{math} et \code{random}. Les fonctions \code{sorted} et \code{sort} sont
autorisées et l'on admettra qu'elles trient une liste de $m$ éléments en $O(m\log m)$.
\item On admet que les opérations élémentaires sur un dictionnaire (lecture, écriture,
test d'appartenance) s'effectuent en $O(1)$ en moyenne.
\item Toutes les complexités demandées sont des complexités \emph{en temps} dans le pire
des cas, sauf indication explicite.
\end{itemize}

\begin{center}
\begin{tikzpicture}[>=Stealth, every node/.style={draw, rounded corners, minimum width=8mm,
 minimum height=5.5mm, inner sep=2pt, font=\small}, level distance=13mm, sibling distance=17mm]
\node (r) {$r{=}0$}
  child {node (c1) {$1$}
    child {node (c2) {$2$}
      child {node (c3) {$3$}}}
    child {node (c4) {$4$}
      child {node (c5) {$5$}}}}
  child {node (c6) {$6$}
    child {node (c7) {$7$}}};
\node[draw=none, right=22mm of r, font=\footnotesize, align=left]
 {\code{parent = [-1, 0, 1, 2, 1, 4, 0, 6]}\\[2pt]
  $h = [0,\,1,\,2,\,3,\,2,\,3,\,1,\,2]$};
\end{tikzpicture}

\emph{Figure 1 --- L'arbre de points de contrôle $A_0$ utilisé comme exemple dans tout le
sujet.}
\end{center}

\bigskip
\noindent\emph{Le sujet comporte 5 parties et 54 questions. Il est vivement conseillé de
traiter les parties dans l'ordre, mais elles restent largement indépendantes.}

\newpage

%=======================================================================
\section*{Partie I --- Points de contrôle, dépôts et fuite d'inactivité}
%=======================================================================

\subsection*{I.1 --- Hauteurs}

Un bloc de numéro $n$ (entier naturel) est un point de contrôle si et seulement si $n$ est
un multiple de $100$~; sa hauteur de point de contrôle vaut alors $n/100$.

\Q Écrire deux fonctions \texttt{est\_checkpoint(n)} et \texttt{hauteur\_bloc(n)} qui prennent en
argument un entier naturel \texttt{n} et renvoient respectivement le booléen indiquant si le
bloc $n$ est un point de contrôle, et l'entier $n/100$ (on supposera dans le second cas que
$n$ est bien un multiple de $100$). Préciser leur complexité.

\Q Écrire une fonction \texttt{hauteur(parent, c)} qui calcule $h(c)$ en remontant les liens
de parenté. Exhiber un \textbf{variant de boucle} entier et positif, et en déduire la
terminaison de la fonction. Donner sa complexité en fonction de $h(c)$.

\Q On souhaite calculer d'un seul coup les hauteurs de \emph{tous} les points de contrôle.
Écrire une fonction \texttt{toutes\_hauteurs(parent)} qui renvoie la liste \texttt{h} telle que
\texttt{h[c]} vaut $h(c)$, en une \textbf{unique} boucle de $c=0$ à $n-1$. Justifier la
correction de l'algorithme en énonçant un invariant de boucle, et préciser à quel endroit
l'hypothèse $\texttt{parent[c]}<c$ est indispensable. Donner la complexité.

\Q Vérifier sur l'arbre $A_0$ de la Figure~1 que \texttt{toutes\_hauteurs} renvoie bien la
liste $h$ indiquée.

\subsection*{I.2 --- Dépôts et supermajorités}

\Q Écrire une fonction \texttt{poids(depots, E)} qui renvoie $w(E)$, puis une fonction
\texttt{est\_supermajorite(depots, E, T)} qui renvoie \texttt{True} si et seulement si $E$ est
une supermajorité. On exige que le test soit effectué \textbf{sans aucun calcul en virgule
flottante}~; justifier l'équivalence utilisée.

\Q On considère les dépôts $\texttt{depots} = \{a:30,\ b:30,\ c:20,\ d:20\}$. Déterminer
$T$, puis dire lesquels des ensembles $\{a,b\}$, $\{a,b,c\}$, $\{a,c,d\}$ sont des
supermajorités.

\Q \textbf{(Lemme d'intersection.)} Soient $E_1$ et $E_2$ deux supermajorités. Démontrer
que \[ w(E_1\cap E_2)\ \ge\ \frac{T}{3}. \] \emph{Indication~:} majorer $w(E_1\cup E_2)$
par $T$ et utiliser $w(E_1\cup E_2)=w(E_1)+w(E_2)-w(E_1\cap E_2)$. Ce lemme est la clé de
toute la partie IV~; on veillera à rédiger la démonstration complètement.

\subsection*{I.3 --- La fuite d'inactivité}

Lorsque plus d'un tiers des validateurs cessent de voter, plus aucun point de contrôle ne
peut être finalisé. Casper prévoit alors une \textbf{fuite d'inactivité}~: à chaque époque,
un validateur qui ne vote pas et dont le dépôt vaut $D$ perd la fraction $p\,D$, où
$p\in\,]0,1[$ est un réel fixé. On note $D_k$ le dépôt d'un validateur inactif après $k$
époques, avec $D_0=D>0$.

\Q Exprimer $D_k$ en fonction de $D$, $p$ et $k$. En déduire la valeur de la somme
$\displaystyle\sum_{k=0}^{+\infty} p\,D_k$ et interpréter ce résultat.

\Q Soit $S$ un réel tel que $0<S<D$. Démontrer que l'ensemble des entiers $k$ tels que
$D_k\le S$ est non vide et que son plus petit élément vaut
\[ k^{\star} \;=\; \left\lceil \frac{\ln(S/D)}{\ln(1-p)} \right\rceil . \]

\Q Écrire une fonction \code{nb_epoques(D, p, S)} qui renvoie $k^{\star}$ \textbf{sans
utiliser de logarithme}, par une simple boucle. Exhiber un variant prouvant la terminaison,
puis montrer que la complexité est un $O\!\left(\log(D/S)\right)$.

\Q Application numérique~: calculer $k^{\star}$ pour $D=1000$, $p=0{,}05$ et $S=500$.

%=======================================================================
\section*{Partie II --- Arbre de points de contrôle et sérialisation des votes}
%=======================================================================

\subsection*{II.1 --- Ancêtres, conflits, plus proche ancêtre commun}

\Q Écrire une fonction \code{est_ancetre(parent, a, b)} qui renvoie \code{True} si et
seulement si \code{a} est un ancêtre de \code{b}. Exhiber un variant de boucle et donner la
complexité en fonction de $h(b)$.

\Q En déduire une fonction \code{conflit(parent, a, b)} qui teste si $a$ et $b$ sont en
conflit. Sur l'arbre $A_0$ de la Figure~1, déterminer si les paires $(2,4)$, $(3,5)$ et
$(1,3)$ sont en conflit.

\Q Écrire une fonction \code{ancetre_commun(parent, a, b, h)}, où \code{h} est la liste des
hauteurs, qui renvoie le \textbf{plus proche ancêtre commun} de $a$ et $b$, c'est-à-dire
l'ancêtre commun de plus grande hauteur. L'algorithme devra procéder en deux temps~: égaliser
les hauteurs, puis remonter simultanément. Donner sa complexité.

\Q Démontrer la correction de \code{ancetre_commun} en énonçant l'invariant de la troisième
boucle. On pourra admettre le résultat suivant, que l'on énoncera précisément~: si $x$ et
$y$ ont même hauteur et sont tous deux ancêtres d'un même n\oe ud, alors $x=y$.

\Q Démontrer l'équivalence~: $a$ et $b$ sont en conflit si et seulement si
$\texttt{ancetre\_commun}(a,b)\notin\{a,b\}$.

\Q Écrire une fonction \code{vote_valide(parent, depots, vote)} qui teste la validité d'un
vote \code{(v, s, t)} au sens des notations.

\subsection*{II.2 --- Sérialisation à longueur variable des hauteurs}

Un vote $\langle \nu, s, t, h(s), h(t)\rangle$ est diffusé sur le réseau~; il faut donc le
\emph{sérialiser}, c'est-à-dire le transformer en une suite d'octets. Comme les hauteurs
sont le plus souvent petites, on n'utilise pas un codage sur un nombre fixe d'octets, mais
le codage \emph{à longueur variable} suivant, dit codage en base $128$.

\begin{defi}
Soit $n\in\mathbb{N}$ et soit $n=\sum_{i=0}^{\ell-1} a_i\,128^{\,i}$ son écriture en base
$128$, avec $0\le a_i\le 127$, $a_{\ell-1}\ne 0$ si $n\ge 1$, et $\ell=1$, $a_0=0$ si
$n=0$. Le \textbf{codage} de $n$ est la liste d'octets
\[ \texttt{encode}(n) \;=\; [\,a_0+128,\ a_1+128,\ \dots,\ a_{\ell-2}+128,\ a_{\ell-1}\,]. \]
Autrement dit, chaque octet porte $7$ bits de poids faible de $n$, et son bit de poids fort
vaut $1$ pour tous les octets sauf le dernier.
\end{defi}

\Q Écrire une fonction \code{encode(n)} conforme à cette définition. Donner
\code{encode(0)}, \code{encode(127)}, \code{encode(128)} et \code{encode(300)}.

\Q Montrer que le nombre de tours de la boucle de \code{encode} vaut $1$ si $n=0$ et
$\lfloor \log_{128} n\rfloor + 1$ si $n\ge 1$. Exhiber un variant de boucle. En déduire que
la longueur de \code{encode(n)} est un $O(\log n)$.

\Q Écrire la fonction réciproque \code{decode(octets)}, de complexité linéaire en la
longueur de la liste.

\Q Démontrer, par récurrence sur $\ell$ (le nombre d'octets produits), que pour tout
$n\in\mathbb{N}$ on a $\texttt{decode}(\texttt{encode}(n)) = n$. On distinguera soigneusement le
cas $n<128$ du cas général.

\Q Le sujet suppose des hauteurs pouvant aller jusqu'à $10^{9}$. Combien d'octets au
maximum le codage d'une hauteur consomme-t-il~? Comparer à un codage sur $8$ octets et
commenter l'intérêt du codage à longueur variable pour un protocole où des milliers de votes
sont diffusés à chaque époque.

%=======================================================================
\section*{Partie III --- Justification et finalisation}
%=======================================================================

On rappelle les définitions du protocole.

\begin{defi}[Casper]\leavevmode
\begin{itemize}[leftmargin=1.4em,itemsep=1pt]
\item Un couple ordonné $(s,t)$ de points de contrôle est un \textbf{lien de supermajorité},
noté $s\to t$, lorsque l'ensemble des validateurs ayant publié un vote valide de source $s$
et de cible $t$ forme une supermajorité. Un lien peut « sauter » des points de contrôle~:
$h(t)>h(s)+1$ est autorisé.
\item Un point de contrôle $c$ est \textbf{justifié} si $c=r$, ou s'il existe un point de
contrôle justifié $c'$ tel que $c'\to c$.
\item Un point de contrôle $c$ est \textbf{finalisé} si $c=r$, ou si $c$ est justifié et
s'il existe un lien $c\to c'$ où $c'$ est un fils \emph{direct} de $c$ (donc
$h(c')=h(c)+1$).
\end{itemize}
\end{defi}

\Q Écrire une fonction \code{liens_supermajorite(votes, depots)} qui prend une liste de
votes valides (éventuellement avec des doublons, un même validateur pouvant émettre
plusieurs fois le même vote) et renvoie l'\textbf{ensemble} des couples $(s,t)$ formant un
lien de supermajorité. On veillera à ne compter le dépôt d'un validateur qu'une seule fois
par couple $(s,t)$. Donner la complexité en fonction du nombre $m$ de votes.

\Q Écrire une fonction \code{enfants_de(parent)} qui renvoie la liste \code{E} telle que
\code{E[c]} soit la liste des fils de $c$. Complexité~?

\Q Écrire une fonction \textbf{récursive} \code{sous_arbre(E, c)} renvoyant la liste des
descendants de $c$ (lui compris). Démontrer que le nombre total d'appels engendrés par
\code{sous_arbre(E, 0)} vaut exactement $n$. Montrer ensuite, sur l'exemple d'un arbre réduit
à un chemin, que le coût des concaténations peut être \textbf{quadratique} en $n$~; proposer
alors une variante \code{sous_arbre_acc} utilisant un accumulateur, de complexité $O(n)$.

\Q Écrire une fonction \code{justifies(parent, liens)} renvoyant l'ensemble des points de
contrôle justifiés. On utilisera un parcours en profondeur itératif à partir de la racine
dans le graphe orienté dont les arcs sont les liens de supermajorité. Donner la complexité
en fonction de $n$ et du nombre $|L|$ de liens.

\Q Démontrer la correction de \code{justifies}~: montrer par double inclusion que
l'ensemble renvoyé est exactement l'ensemble des points de contrôle justifiés. Pour
l'inclusion réciproque, on raisonnera par récurrence sur la longueur d'une chaîne de
justification menant de $r$ à $c$.

\Q Écrire une fonction \code{finalises(parent, liens, J, h)} renvoyant l'ensemble des
points de contrôle finalisés, où \code{J} est l'ensemble des justifiés. Donner la
complexité.

\Q Démontrer que si $c$ est finalisé et $c\ne r$, alors le fils direct $c'$ intervenant
dans la définition est lui-même justifié.

\Q On considère l'arbre $A_0$ de la Figure~1, les dépôts de la question 6, et la liste de
votes où chacun des validateurs $a$, $b$, $c$ émet les trois votes $(0,1)$, $(1,2)$ et
$(2,3)$. Déterminer l'ensemble des liens de supermajorité, l'ensemble des justifiés, puis
l'ensemble des finalisés. Justifier chaque appartenance.

\subsection*{III.4 --- La règle de choix de fourche}

La règle de choix de fourche de Casper s'énonce~: \emph{suivre la chaîne contenant le point
de contrôle justifié de plus grande hauteur.}

\Q Écrire une fonction \code{choix_fourche(J, h)} renvoyant le point de contrôle justifié de
plus grande hauteur. Complexité~?

\Q On admet la propriété (iv) démontrée en partie IV~: \emph{tant que moins d'un tiers des
validateurs (en dépôt) n'enfreint aucun commandement, il existe au plus un point de contrôle
justifié de hauteur donnée.}
\begin{enumerate}[label=(\alph*),leftmargin=2em,itemsep=1pt]
\item Démontrer que tout point de contrôle justifié $c$ est relié à la racine par une chaîne
$r=b_0\to b_1\to\dots\to b_k=c$ de liens de supermajorité dont \emph{tous} les termes sont
justifiés~; en déduire que chaque $b_i$ est un ancêtre de $b_{i+1}$ et que la suite
$\big(h(b_i)\big)_{0\le i\le k}$ est strictement croissante.
\item En déduire que le maximum calculé par \code{choix_fourche} est atteint en un
\textbf{unique} n\oe ud, et que la règle de choix de fourche est donc bien définie.
\end{enumerate}

%=======================================================================
\section*{Partie IV --- Détection des infractions et sûreté responsable}
%=======================================================================

\begin{defi}[Les deux commandements de Casper]
Un validateur $\nu$ ne doit publier deux votes distincts $\langle\nu,s_1,t_1\rangle$ et
$\langle\nu,s_2,t_2\rangle$ vérifiant l'une des deux conditions suivantes~:
\begin{itemize}[leftmargin=2.2em,itemsep=1pt]
\item[\textbf{(I)}] $h(t_1)=h(t_2)$ \quad (deux votes distincts pour une même hauteur de
cible)~;
\item[\textbf{(II)}] $h(s_1)<h(s_2)<h(t_2)<h(t_1)$ \quad (un vote strictement emboîté dans
un autre).
\end{itemize}
Tout validateur enfreignant l'un de ces commandements perd la totalité de son dépôt.
\end{defi}

\subsection*{IV.1 --- Détection algorithmique}

\Q Écrire une fonction \code{infraction_I(votes, h)} qui renvoie un témoin d'infraction au
commandement (I) — c'est-à-dire un triplet $(\nu, t_1, t_2)$ avec $t_1\ne t_2$ et
$h(t_1)=h(t_2)$ tous deux ciblés par $\nu$ — ou \code{None} s'il n'y en a pas. La complexité
devra être en $O(m)$ en moyenne, où $m$ est le nombre de votes. Justifier.

\Q Écrire une fonction \code{infraction_II_naif(votes, h)} détectant une infraction au
commandement (II) par double boucle. Donner sa complexité.

\Q On veut faire mieux. Soit un validateur $\nu$ et la liste $L$ de ses votes, triée par
hauteur de source croissante. Démontrer l'équivalence suivante~: \emph{$\nu$ enfreint (II)
si et seulement s'il existe un vote $(s_2,t_2)$ de $L$ tel que
\[ \max\{\,h(t_1)\ :\ (s_1,t_1)\in L,\ h(s_1)<h(s_2)\,\}\ >\ h(t_2). \]}
On prendra garde au fait que la condition $h(s_2)<h(t_2)$ est automatiquement vérifiée pour
un vote valide.

\Q En déduire une fonction \code{infraction_II(votes, h)} de complexité $O(m\log m)$,
utilisant un tri puis un unique parcours maintenant le maximum courant des $h(t)$ déjà
rencontrés. On prendra soin de traiter les votes de \emph{même} hauteur de source par
groupes, afin de respecter l'inégalité \emph{stricte} $h(s_1)<h(s_2)$. Énoncer l'invariant
de la boucle principale et justifier la complexité.

\Q Sur l'arbre $A_0$, un validateur $\nu$ publie les deux votes $(0,3)$ et $(1,2)$. Quel
commandement enfreint-il~? Même question pour les votes $(0,2)$ et $(0,4)$.

\subsection*{IV.2 --- Sûreté responsable}

On suppose désormais que \textbf{moins d'un tiers} des validateurs (en dépôt) enfreint un
commandement.

\Q Démontrer les propriétés suivantes.
\begin{itemize}[leftmargin=2.2em,itemsep=1pt]
\item[(i)] Si $s_1\to t_1$ et $s_2\to t_2$ sont deux liens de supermajorité distincts, alors
$h(t_1)\ne h(t_2)$.
\item[(ii)] Si $s_1\to t_1$ et $s_2\to t_2$ sont deux liens de supermajorité distincts, alors
l'inégalité $h(s_1)<h(s_2)<h(t_2)<h(t_1)$ est impossible.
\end{itemize}
\emph{Indication~:} appliquer le lemme d'intersection de la question 7 aux deux ensembles de
validateurs soutenant les deux liens.

\Q En déduire~:
\begin{itemize}[leftmargin=2.2em,itemsep=1pt]
\item[(iii)] pour toute hauteur $n$, il existe au plus un lien de supermajorité $s\to t$ avec
$h(t)=n$~;
\item[(iv)] pour toute hauteur $n$, il existe au plus un point de contrôle justifié de
hauteur $n$.
\end{itemize}

\Q \textbf{(Théorème de sûreté responsable.)} Soient $a_m$ et $b_n$ deux points de contrôle
finalisés \emph{en conflit}, de fils directs justifiés respectifs $a_{m+1}$ et $b_{n+1}$.
Quitte à échanger les rôles, on suppose $h(a_m)<h(b_n)$. En considérant la chaîne de
justification $r\to b_{i_1}\to b_{i_2}\to\dots\to b_n$ et en introduisant le plus petit
indice $j$ tel que $h(b_j)>h(a_{m+1})$, démontrer que l'on aboutit à une contradiction avec
(ii) ou (I). Conclure~: \emph{deux points de contrôle en conflit ne peuvent pas être tous
deux finalisés sans qu'au moins un tiers du dépôt total ne soit détruit.}

\Q Pourquoi le cas $h(a_m)=h(b_n)$ se traite-t-il immédiatement~? Quel commandement est
alors enfreint~?

\subsection*{IV.3 --- Recherche probabiliste d'un validateur fautif}

Un client léger ne peut pas vérifier tous les validateurs. Il tire donc au hasard,
\textbf{avec remise et de façon indépendante}, des validateurs uniformément parmi les $N$
validateurs, et vérifie pour chacun s'il est fautif. On suppose qu'une proportion
$q\in\,]0,1[$ des validateurs est fautive. On note $X$ le rang du premier tirage amenant un
validateur fautif.

\Q Justifier que $X$ suit une loi géométrique de paramètre $q$. Donner $P(X=k)$ pour
$k\ge 1$, puis $P(X>k)$, et enfin $\mathbb{E}[X]$.

\Q En déduire que l'algorithme de recherche s'arrête \textbf{presque sûrement}, c'est-à-dire
que $P(X=+\infty)=0$.

\Q Écrire une fonction \code{recherche_fautif(validateurs, est_fautif)} qui implémente cette
recherche à l'aide de \code{random.randrange}, où \code{est_fautif} est une fonction
booléenne. Que vaut l'espérance du nombre d'appels à \code{est_fautif} lorsque exactement un
tiers des validateurs est fautif~? Pourquoi ne peut-on pas majorer le nombre de tours dans
le pire des cas~?

%=======================================================================
\section*{Partie V --- Coût minimal d'une attaque (programmation dynamique)}
%=======================================================================

D'après la question 7, une attaque réussie contre la sûreté de Casper coûte à l'attaquant au
moins $T/3$. Un attaquant rationnel cherche donc à \textbf{minimiser} le dépôt qu'il met en
jeu tout en réunissant une supermajorité. Cela conduit au problème d'optimisation suivant.

\begin{defi}[\textsc{Coût-Attaque}]
\emph{Donnée~:} une liste $d=[d_0,\dots,d_{n-1}]$ d'entiers strictement positifs (les
dépôts) et un entier $C\ge 0$ (le seuil, en pratique $C=\lceil 2T/3\rceil$).\\
\emph{Question~:} déterminer
$\displaystyle \mu(d,C)=\min\Big\{\textstyle\sum_{i\in S} d_i \ :\ S\subseteq
\{0,\dots,n-1\},\ \sum_{i\in S} d_i\ \ge\ C\Big\}$, avec $\mu=+\infty$ si aucun $S$ ne
convient.
\end{defi}

\Q Justifier que $C=\lceil 2T/3\rceil$ est bien le plus petit entier $C$ tel que
$w(E)\ge C \iff E$ est une supermajorité. Écrire l'expression Python de $C$ à partir de
\code{T} en n'utilisant que des divisions entières.

\Q Le problème de \textbf{décision} associé est~: « existe-t-il $S$ avec
$\sum_{i\in S} d_i = C$~? » (problème \textsc{Somme-de-Sous-Ensemble}). Démontrer que ce
problème appartient à la classe \textbf{NP} en explicitant un certificat et un vérificateur
polynomial. On rappelle (et on admettra) que \textsc{Somme-de-Sous-Ensemble} est
NP-complet~; que peut-on en déduire quant à l'existence d'un algorithme polynomial pour
\textsc{Coût-Attaque}, sous l'hypothèse $\mathrm{P}\ne\mathrm{NP}$~?

\Q On pose, pour $0\le i\le n$ et $0\le j\le C$~:
\[ m_i(j) \;=\; \min\Big\{\textstyle\sum_{k\in S} d_k \ :\ S\subseteq\{0,\dots,i-1\},\
\min\big(C,\textstyle\sum_{k\in S} d_k\big) = j \Big\}\ \in\ \mathbb{N}\cup\{+\infty\}. \]
Autrement dit, $m_i(j)$ est la plus petite somme réalisable avec les $i$ premiers dépôts
dont la valeur \emph{écrêtée à $C$} vaut exactement $j$. Donner $m_0$, puis établir la
relation de récurrence exprimant $m_{i+1}$ en fonction de $m_i$. Expliquer pourquoi
l'écrêtage $\min(C,\cdot)$ est nécessaire pour borner la taille du tableau, et pourquoi il
ne fait perdre aucune solution optimale.

\Q Exprimer $\mu(d,C)$ à l'aide de $m_n$.

\Q Écrire une fonction \code{cout_min_attaque(d, C)} implémentant cette récurrence avec un
\textbf{unique} tableau de taille $C+1$ mis à jour en place. Justifier précisément le
\textbf{sens de parcours} de l'indice $j$ et montrer sur un exemple ce qui se produirait si
l'on parcourait dans l'autre sens.

\Q Démontrer la correction de l'algorithme en énonçant et en prouvant, par récurrence sur
le nombre $i$ de dépôts déjà traités, l'\textbf{invariant} vérifié par le tableau à la fin
de la $i$-ème itération de la boucle externe. Justifier également la terminaison.

\Q Donner la complexité en temps et en espace de \code{cout_min_attaque}. Expliquer
pourquoi cet algorithme est dit \textbf{pseudo-polynomial} et pourquoi cela ne contredit pas
la question 46 (NP-complétude). Que devient la complexité si les dépôts sont exprimés en
\emph{wei}, la plus petite unité de l'ether, soit $10^{18}$ wei par ether~?

\Q Application numérique~: on prend $d=[12,\,7,\,20,\,5,\,16]$. Calculer $T$, puis $C$, puis
$\mu(d,C)$ en explicitant un sous-ensemble optimal et en détaillant le contenu du tableau à
la fin de chaque itération pour la version réduite $d=[3,5,7]$, $C=8$.

\Q Comparer $\mu(d,C)$ au seuil $C$ sur cet exemple. Un attaquant peut-il toujours
mobiliser \emph{exactement} $C$~? Proposer un jeu de dépôts pour lequel $\mu(d,C)>C$
strictement, et le justifier. Rapprocher enfin ce résultat de la question 7~: quelle part du
dépôt mobilisé est effectivement \emph{détruite} lors d'une violation de la sûreté~?

\Q \textbf{Question de synthèse (à traiter en quelques lignes).} Casper garantit la
\emph{sûreté responsable} mais pas la \emph{vivacité} en toutes circonstances. À la lumière
des parties III, IV et V, expliquer en quoi le mécanisme de fuite d'inactivité de la partie I
restaure la vivacité, et pourquoi il fragilise en retour la garantie de sûreté démontrée à
la question 49.

\vfill
\begin{center}\small\textsc{Fin de l'épreuve}\end{center}
\end{document}
